%% =====================================================================
%% One page on the proof of concept, plus an appendix page.
%% AD Assurance Lab, Western Michigan University.
%% =====================================================================
\documentclass[10pt,twocolumn]{article}
\usepackage[letterpaper,margin=0.72in]{geometry}
\usepackage[T1]{fontenc}
\usepackage{mathptmx}
\usepackage{graphicx}
\usepackage{amsmath}
\usepackage{xcolor}
\usepackage{url}
\usepackage[hidelinks]{hyperref}
\usepackage{caption}
\usepackage{subcaption}
\usepackage{authblk}
\usepackage{booktabs}
\usepackage{titlesec}
\titlespacing*{\section}{0pt}{1.6ex plus .2ex minus .2ex}{1.0ex plus .1ex}
\titlespacing*{\subsection}{0pt}{1.3ex plus .2ex minus .2ex}{0.7ex plus .1ex}
\captionsetup{font=small,labelfont=bf}
\renewcommand\Authfont{\normalsize}
\renewcommand\Affilfont{\small\itshape}
\setlength{\affilsep}{0.3em}
\pagestyle{plain}
\renewcommand{\topfraction}{0.9}
\renewcommand{\bottomfraction}{0.9}
\renewcommand{\textfraction}{0.07}
\renewcommand{\floatpagefraction}{0.75}
\setcounter{topnumber}{3}
\setcounter{totalnumber}{5}
\raggedbottom

\title{\bfseries\large Rebuilt Roads Move a Driving Policy as Much as Fog Does:
A First Acceptance Measurement on Real Frames}

\author[1]{Zachary D. Asher}
\affil[1]{AD Assurance Lab, Western Michigan University, Kalamazoo, MI, USA}
\date{\vspace{-1.0em}\small September 20, 2026}

\begin{document}
\maketitle

\begin{abstract}
\noindent
Our verified disturbance family needs a clear frame and a degraded frame at one
camera pose, and no published dataset supplies that pair. Neural reconstruction
can, so we built the acceptance test first and measured the reconstruction with
it. On a real recorded clip, the rebuilt frame moves an open reasoning policy's
planned path by 0.25\,m at a three-second horizon. A deliberately heavy fog on
the real frame moves it by 0.33\,m. The substrate error is the size of the
disturbance we wanted to certify, and at six seconds its worst pose is larger. The policy's own sampling moves it further still, by 1.67\,m.
\end{abstract}

\section{Why this measurement}
The lab verifies a policy over a one-parameter family between two images of one
scene. That needs $x_\text{clear}$ and $x_\text{cond}$ at an identical camera
pose, which no published adverse-weather dataset supplies. Reconstruction
promises those endpoints from any logged drive, and it carries a trap we have
measured once. An analytic fog model reproduced simulator fog at road-ROI
$R^2=0.848$, a good image score. It then drove one policy 23.8 times harder than
real fog, and stayed faithful for another. Image fidelity is not the property
that matters, so the acceptance test is behavioral and it comes first.

That test has three steps, and this note reports the middle one. We first
measure the spread between two real frames, and then check the rebuilt clear
frame against the real one. A rebuilt adverse frame comes last, and only after
those two pass.

\section{What we ran}
\textbf{Data.} One clip from the public NVIDIA reconstruction collection, which
ships the real forward video and the recorded rig poses beside the
reconstruction. We used its first eight seconds at 10\,Hz and 1920 by 1080, from a night
drive at about 18\,m/s.

\textbf{Rendering.} The 2026 engine renders along the recorded trajectory, one
image per camera frame, so rendered frame $N$ sits at the pose of recorded frame
$N$. That is the pairing the test needs.

\textbf{Policy.} Alpamayo 1.5, an open 10-billion-parameter reasoning model that
predicts 6.4\,s of trajectory and a short reasoning trace. On the real frames it
predicts 19.6\,m of travel in one second, which matches the recorded speed. It
needs 22.5\,GB of our 32\,GB card and answers in about 0.5\,s per pose.

\textbf{Control.} The known-bad disturbance from the protocol, an analytic
Koschmieder fog on the real frame, which exists so the instrument has something
it must reject.

\section{What we measured}
We run the policy over the same 65 poses on each image set. The configuration
is repeatable: greedy text, one trajectory, one fixed seed. A repeat run
reproduces the first exactly, so the floor is 0.00\,m. Every centimeter of
difference then belongs to the image.

Fig.~\ref{fig:gap} and Table~\ref{tab:main} give the result. The rebuilt frame
moves the planned path by 0.25\,m at three seconds, against 0.33\,m for heavy
fog on the real frame. At six seconds the mean is smaller, 0.80\,m against
1.20\,m, but the worst rebuilt pose is larger, 4.20\,m against 3.20\,m.

\begin{figure}[t]
  \centering
  \includegraphics[width=\columnwidth]{figures/gap_horizon.pdf}
  \caption{Three ways to move the same policy at the same poses. The
  reconstruction sits beside the disturbance it has to carry, and the triangles
  mark the worst pose of each.}
  \label{fig:gap}
\end{figure}

\section{What follows from it}
\textbf{This reconstruction cannot yet carry this disturbance.} A certificate
over a fog-sized family, computed on this rebuild, would rest on a substrate
whose own error is the size of the family. The clear-frame step exists to catch
exactly that, and it caught it on the first clip we tried.

\textbf{It is a measurement, not a verdict.} One clip, one camera, one policy,
and the engine has settings we have not explored. The number to improve is now
defined, and it is behavioral rather than photometric.

\textbf{Hold the sampling still, or measure nothing.} With sampling on, the
policy disagrees with itself by 1.67\,m at three seconds, which beats either
image effect. Our protocol requires the repeatable configuration, and records
that spread beside every gap.

\clearpage
\appendix
\section{Method, defects and limits}

\begin{figure}[h]
  \centering
  \begin{subfigure}{0.49\columnwidth}
    \includegraphics[width=\linewidth]{figures/frame_clear.png}
    \caption{recorded}
  \end{subfigure}\hfill
  \begin{subfigure}{0.49\columnwidth}
    \includegraphics[width=\linewidth]{figures/frame_rebuilt.png}
    \caption{rebuilt}
  \end{subfigure}\\[3pt]
  \begin{subfigure}{0.49\columnwidth}
    \includegraphics[width=\linewidth]{figures/frame_fog.png}
    \caption{the known-bad fog}
  \end{subfigure}
  \caption{One pose, three image sources. The policy sees each in turn, and the
  instrument compares what it plans.}
  \label{fig:frames}
\end{figure}

\begin{table}[h]
\centering
\small
\begin{tabular}{lrrrr}
\toprule
& \multicolumn{2}{c}{Rebuilt frame} & \multicolumn{2}{c}{Known-bad fog} \\
\cmidrule(lr){2-3}\cmidrule(lr){4-5}
Horizon & mean & worst & mean & worst \\
\midrule
1\,s & 0.04\,m & 0.25\,m & 0.05\,m & 0.34\,m \\
3\,s & 0.25\,m & 1.29\,m & 0.33\,m & 1.52\,m \\
6\,s & 0.80\,m & 4.20\,m & 1.20\,m & 3.20\,m \\
\bottomrule
\end{tabular}
\caption{Sixty-five poses of one clip, against the real frame at the same pose.
The policy's own sampling spread is 0.23\,m, 1.67\,m and 4.37\,m.}
\label{tab:main}
\end{table}

\subsection{How to repeat it}
The scripts live in \texttt{scripts/} of the reconstruction repository.
\begin{verbatim}
render --artifact-path S --camera-id C
       --image-scale 1.0 --frame-step 3

make_rendered_frames.py --rendered R
       --usdz S --rendered-step 3

run_policy.py  --frames F --samples 1
       --temperature 0.0 --stride 1

response_gap.py --a REAL --b REBUILT
       --horizon 3.0
\end{verbatim}

\subsection{A defect the instrument caught}
The first run predicted a path that ran backwards. The scene file stores the rig
pose in the world, and we had inverted it, which mirrors the ego history. The
policy followed that mirrored history faithfully. No image score would have
shown this, and the behavioral check showed it in one plot.

\subsection{Limits, stated plainly}
One clip. One camera of the four the policy expects. Eight seconds. Night, clear
weather, one road. The real frames come from a compressed video, and the rebuilt
frames come from the engine's default renderer at native resolution. The fog is
our control, not a measured disturbance. Four samples give a coarse spread. None
of this is a certificate.

\begin{figure}[b]
  \centering
  \includegraphics[width=\columnwidth]{figures/gap_time.pdf}
  \caption{The reconstruction gap at a three-second horizon, at every pose. It
  never reaches zero, and it peaks where the road bends.}
  \label{fig:time}
\end{figure}

\end{document}
